\documentclass[12pt]{article}
\usepackage[T1]{fontenc}
\usepackage[utf8]{inputenc}
\usepackage{lmodern}
\usepackage{textcomp}
\usepackage{enumitem}
\usepackage{geometry}
\geometry{margin=1in}
\begin{document}
\begin{center}
Answer Key - Constitutional Law - Graduate Level Assessment 2 \
\small (Answers are ordered exactly as in the question paper.)
\end{center}

\section*{Multiple Choice}
\begin{enumerate}[label=\arabic*.]
\item \textbf{Answer:} A
\\ \textit{Options:} A) An implied-in-fact contract.; B) An implied-in-law contract.; C) An express contract.; D) No contract.
\\ \textit{Note:} The correct answer is an implied-in-fact contract because the acceptance of legal services by the client, despite the absence of a formal fee agreement, establishes an understanding that the client will compensate the attorney for those services based on the reasonable value of the work performed.
\item \textbf{Answer:} A
\\ \textit{Options:} A) Cold pursuit; B) Incident to a lawful arrest; C) Plain view; D) Stop and frisk
\\ \textit{Note:} Cold pursuit is not recognized as a warrantless search exception because it refers to the immediate need to apprehend a suspect, which does not inherently justify a search without a warrant under constitutional law.
\item \textbf{Answer:} D
\\ \textit{Options:} A) Contingent remainders; B) Executory interests; C) Vested remainders subject to open; D) Reversions
\\ \textit{Note:} The correct answer, "Reversions," pertains to the concept in property law where ownership of a property reverts back to the grantor upon the occurrence of a specified event, such as the intent to commit a felony, which aligns with the understanding of conditional interests in property rights.
\item \textbf{Answer:} C
\\ \textit{Options:} A) The proceedings are conducted in secret.; B) There is no right to counsel.; C) There is a right to Miranda warnings.; D) There is no right to have evidence excluded.
\\ \textit{Note:} The correct answer is based on the fact that Miranda warnings are not required during grand jury proceedings because these proceedings are not considered custodial interrogation under the Fifth Amendment, which only applies to situations where an individual is formally in custody and subjected to interrogation by law enforcement.
\item \textbf{Answer:} A
\\ \textit{Options:} A) Bill of attainder; B) Ex post facto law; C) Statutes; D) Ordinances
\\ \textit{Note:} A bill of attainder is a legislative act that specifically targets an individual or group, imposing punishment without the benefit of a judicial trial, thereby violating the constitutional principle of due process.
\end{enumerate}

\section*{True / False}
\begin{enumerate}[label=\arabic*.]
\setcounter{enumi}{5}
\item \textbf{Answer:} False
\\ \textit{Options:} A) True; B) False
\\ \textit{Note:} The necessary and proper clause cannot stand alone as a constitutional basis for federal law because it is intended to be used in conjunction with other enumerated powers granted to Congress in order to carry out its legislative functions effectively.
\item \textbf{Answer:} False
\\ \textit{Options:} A) True; B) False
\\ \textit{Note:} An irrevocable written offer, even if consideration is provided, cannot last indefinitely because it is subject to a reasonable time limit for acceptance based on the circumstances of the offer and the nature of the transaction.
\item \textbf{Answer:} False
\\ \textit{Options:} A) True; B) False
\\ \textit{Note:} The statement is false because the reasonable person standard in negligence cases is based on the conduct expected of an average person in similar circumstances, and a defendant's IQ alone does not automatically determine their ability to meet this standard, as it also considers other factors such as the situation and individual circumstances.
\item \textbf{Answer:} False
\\ \textit{Options:} A) True; B) False
\\ \textit{Note:} A substantive due process problem specifically relates to the fundamental rights and liberties being infringed upon, rather than the mere differential treatment under the law, which may instead implicate equal protection concerns.
\item \textbf{Answer:} True
\\ \textit{Options:} A) True; B) False
\\ \textit{Note:} The best evidence rule is primarily concerned with the reliability and authenticity of primary documents, and it allows for exceptions in cases involving collateral matters, public records, or undisputed contents to promote judicial efficiency and access to evidence.
\end{enumerate}

\section*{Long Form Response}
\begin{enumerate}[label=\arabic*.]
\setcounter{enumi}{10}
\item \textbf{Answer:} The constitutional requirements for the termination of a parent's custody rights and welfare benefits hinge critically on the principles of due process, which necessitate prior notice and an evidentiary hearing. The U.S. Supreme Court has established that parents have a fundamental right to raise their children, and any state action that seeks to sever this relationship must be accompanied by adequate procedural safeguards. Prior notice ensures that parents are informed of the actions being taken against them, allowing them the opportunity to prepare a defense. Additionally, an evidentiary hearing provides a platform for the presentation of evidence and arguments, ensuring that the decision to terminate custody rights is not made arbitrarily but instead based on a thorough examination of the facts. Thus, both prior notice and hearings are essential to uphold the constitutional protections afforded to parents under the Fourteenth Amendment.
\\ \textit{Note:} correct because it emphasizes the due process requirements of prior notice and evidentiary hearings, which are essential for protecting the fundamental parental rights under the Constitution when the state seeks to terminate custody rights or welfare benefits.
\item \textbf{Answer:} In the context of auctions, the legal principle governing bid retraction allows a bidder to withdraw their bid at any point until the auctioneer's hammer falls, signifying the conclusion of the bidding process. This principle is crucial, particularly during situations such as a competitor's liquidation, where the financial instability of competing bidders may influence the market dynamics and bidding strategies. For bidders like Fernandez, this flexibility ensures that they can reassess their commitment in light of new information or changing circumstances, ultimately fostering a more equitable bidding environment. However, it also raises concerns about the potential for strategic retractions that could disrupt the auction's integrity and fairness. Understanding these implications is essential for bidders to navigate legal and ethical boundaries in competitive auction settings.
\\ \textit{Note:} The answer correctly identifies that bidders can retract their bids until the auction concludes, which is especially relevant in liquidations where financial instability influences bidding strategies, thus providing bidders like Fernandez the necessary flexibility to evaluate their positions.
\end{enumerate}

\vfill
\noindent\textit{End of Answer Key}
\end{document}